\documentclass[11pt,a4paper]{article}
\usepackage[T1]{fontenc}
\usepackage[utf8]{inputenc}
\usepackage{lmodern}
\usepackage[margin=29mm]{geometry}
\usepackage{amsmath,amssymb,amsthm,booktabs,array,enumitem,microtype}
\usepackage[hidelinks]{hyperref}
\usepackage{xcolor}
\newtheorem{proposition}{Proposition}
\newcommand{\eps}{\varnothing}
\newcommand{\lex}{\leq_{\mathrm{lex}}}
\newcommand{\initial}{\sqsubseteq}
\newcommand{\starred}{W^{\mathcal D}_*}
\newcommand{\lv}{\operatorname{ht}}
\newcommand{\Int}{\operatorname{Int}}
\newcommand{\ms}{\operatorname{ws}}
\newcommand{\sym}[1]{v_{#1}}
\setlength{\parindent}{0pt}
\setlength{\parskip}{0.45em}
\title{An ordered marked-point obstruction in the $Q$-construction\\[3pt]
\large A note on the proof of Lemma~27 of Todorcevic--Tyros}
\author{Verification note for the authors}
\date{10 October 2026}
\begin{document}
\maketitle
\vspace{-1.5em}

\begin{abstract}
We describe a concrete choice of a starred mixed-product point for which the
formula defining the tree component $S_w$ of $Q(w)$, on page~28 of
\emph{A Dual Ramsey Theorem for Trees} (arXiv:2207.14599v1),
violates clause~(ii) of the definition of a skew subtree.
The marked points satisfy the \emph{vector $1$-complete} condition of
Definition~20 when their coordinates are indexed in the paper's
length-first, reverse-lexicographic order. The example works at every
complete-tree height $n\geq3$; hence merely increasing $n$ does not
repair this particular well-definedness assertion. The argument concerns
the displayed $Q$-construction, \emph{not} the conclusion of Lemma~27 or the
main dual Ramsey theorem.
\end{abstract}

\section*{The two relevant conventions}
Let $2^{<h}$ be the binary tree of sequences of length below $h$,
with root $\eps$ and $s\mathbin{^\frown}t$ denoting concatenation.
Write $\initial$ for the prefix relation. The paper's auxiliary order
$\preceq$ (Section~3.1, p.~7) compares lengths first, then uses
\emph{reverse} lexicographic order on equal-length sequences:
\[
  1\prec 0, \qquad |0|=|1|=1.
\]
For a subtree $A$, its intrinsic height at $u\in A$ is
$\lv_A(u)=|\{v\in A:v\sqsubset u\}|$. Clause~(ii) of skewness says
\begin{equation}\label{eq:skew}
  \lv_A(u)=\lv_A(v),\quad u\lex v
  \quad\Longrightarrow\quad |u|\leq |v|.
\end{equation}

Definition~20 (p.~24) imposes an additional restriction on the
$D_2$-marked coordinates of $\starred(b,h,\Lambda)$:
the ordered vector of marked points must belong to
$\mathrm{CT}_1(\mathbf T^{|D_2|}_{b,h})$.
For two coordinates, this amounts to each singleton point being
$1$-complete, with the points' lengths \emph{nondecreasing in
coordinate order}. In particular, equal lengths are \emph{not} required.
This starred condition is essential to the example below.

\section*{A source satisfying the signature hypotheses}
Fix any $n\geq3$ and $N>n$, and let $\Lambda$ be any nonempty
alphabet. Set $b=2$, $l=1$, $m'=0$, and choose any $m\geq2$.
Let
\[
 T=2^{<n}\subseteq 2^{<N}, \qquad t_0=\eps.
\]
This $T$ is an $n$-complete skew tree, and its canonical isomorphism
$I_T:2^{<n}\to T$ is the identity. Define a $T$-variable word $f$ on
$2^{<N}$ by
\[
  f(u)=\sym{u\upharpoonright\min\{|u|,n-1\}}
  \qquad (u\in 2^{<N}).
\]
Every variable $\sym t$ occurs at its own root $t\in T$ and only at
extensions of $t$, so $\ms(f)=T$ and $f\in W_{v,n}(2,N,\Lambda)$.

Now let
\[
  S=\{\eps,0,1\}, \qquad g(u)=\sym\eps
  \quad(u\in 2^{<N}).
\]
The tree $S$ is semi-complete skew with $\Int(S)=\{\eps\}$;
$g$ has that singleton as its variable support. Moreover,
$[g]_\Lambda\subseteq[f]_\Lambda$: every constant evaluation of $g$
is obtained from $f$ by assigning the same letter to all variables.
Thus $(S,g)\in W^*_{v,1}(f)$ and its signature $\mathcal S$
is observable by $f$ at $t_0=\eps$.

For this cut, the strict auxiliary initial segment $D$ of
Definition~26 is empty, and no leaf of $S$ lies outside the cone
of $\eps$. Hence
\[
  S'=\Int(S)=\{\eps\},\qquad g'=g\upharpoonright D
  \text{ is the empty function},\qquad
  R=\{0,1\}.
\]
Here $R$ is exactly the set used on page~28 to select the $D_2$
coordinates. The minimal points of
$T\setminus\{t\in T:t\preceq t_0\}$ are $0$ and $1$.
\emph{In the prescribed auxiliary order}, they are enumerated as
\[
 t_1=1,\qquad t_2=0,\qquad d=2,\quad D_2=\{1,2\},\quad D_0=D_1=\emptyset.
\]

\section*{An admissible starred mixed-product point}
In the proof of Lemma~27, $n'=n-m'=n$. Pick any letter $a\in\Lambda$,
and let $w_1,w_2$ be the constant words with value $a$. Choose the
$D_2$ points
\[
                 s_1=\eps, \qquad s_2=0.
\]
Then
\[
      w=((w_1,w_2),\emptyset,(s_1,s_2))
        \in \starred(2,n',\Lambda).
\]
Indeed, each $\{s_i\}$ is a singleton $1$-complete skew subtree,
and $(\{s_1\},\{s_2\})$ is vector $1$-complete: its only nonvacuous
cross-coordinate length requirement is
$|s_1|=0\leq1=|s_2|$.
The selected point tuple therefore satisfies \emph{Definition~20},
not merely membership in the larger space $W^{\mathcal D}$.

The cone maps of page~28 are $P_i(z)=I_T(I_T^{-1}(t_i)
\mathbin{^\frown}z)$. In our full homogeneous support,
$P_i(z)=t_i\mathbin{^\frown}z$ wherever defined. Both selected tails
are safely in their domains, since $|t_i|+|s_i|<n$:
\[
\begin{array}{c|c|c|c|c}
\text{coordinate }i&t_i&s_i&|s_i|&P_i(s_i)\\\hline
1&1&\eps&0&1\\
2&0&0&1&00
\end{array}
\]
The displayed formula for the tree component of $Q(w)$ gives
\begin{equation}\label{eq:tree}
 S_w=\Int(S')\cup\{t_0\}\cup
       \{P_i(s_i):i\in D_2\}
     =\{\eps,1,00\}.
\end{equation}

\begin{proposition}
The tree $S_w$ in \eqref{eq:tree} is not skew and therefore cannot be
the tree component of any element of $W^*_{v,1}(2,N,\Lambda)$.
\end{proposition}
\begin{proof}
The two nonroot nodes $00$ and $1$ both have the root $\eps$
as their only strict predecessor in $S_w$. Thus
$\lv_{S_w}(00)=\lv_{S_w}(1)=1$. In the usual lexicographic order,
$00<_{\mathrm{lex}}1$, but $|00|=2>1=|1|$, contradicting
\eqref{eq:skew}. No choice of the word component $g_w$ can alter
this tree-level failure.
\end{proof}

All choices exist for every $n\geq3$ and $N>n$, independently of a
numerical lower bound $n_0$. A constant colouring $c$ can be used,
if desired, to satisfy the colouring hypothesis of Lemma~27.
The issue is the assertion immediately after the definition of $Q$
that it is well defined as a map into $[\mathcal S,f]$;
consequently the subsequent pullback $c'(w)=c(Q(w))$ is not defined
on this admissible input as written. The unrestricted formula
for $P_i$ has a separate tail-domain issue near the top level;
the selected $s_1,s_2$ here do \emph{not} use any problematic
out-of-domain tail.

\section*{Scope, verification, and possible repairs}
The example is a counterexample to this particular \emph{proof step},
not to the existence assertion in Lemma~27 or Theorem~25.
It also does not establish that every alternative construction of $Q$
necessarily fails. Two changes seem worth examining:
\begin{enumerate}[leftmargin=1.8em,itemsep=2pt]
 \item Revisit the reverse-lexicographic tie-break in $\preceq$
 (and therefore the frontier indexing). Forward lexicographic
 indexing aligns the vector $1$-complete length inequality with
 the skew clause~(ii) in this minimal example. Such a change needs
 an independent check throughout the paper; other audited uses of
 the printed order are also delicate.
 \item Alternatively, modify the allowed $D_2$ points or the map $Q$
 so that the \emph{projected} terminal nodes obey the required
 lexicographic length monotonicity. A sufficient geometric
 candidate is to place all projected terminals at a common intrinsic
 level of $T$ strictly above the old interior vertices: skewness of
 $T$ then controls their ambient lengths. This restriction must be
 reconciled with the starred mixed-product Ramsey step and its span.
\end{enumerate}
Neither option has yet been proved to repair the full argument.

\paragraph{Formal verification and a question for the authors.}
The Lean development at
\url{https://github.com/janhubicka/lean-dualtree}
checks the complete-height-three instance with the \emph{literal}
marker-to-frontier indexing, not just an analogous three-node tree.
In particular, the theorem
\texttt{sourceQ\_condIIB\_false} in
\texttt{QLengthActualSkewFailure}
proves failure of skew clause~(ii) for the actual projected $Q$ tree.
The theorem \path{actual_marked_points_in_starred_domain}
in \texttt{QLengthStarredDomain} verifies the vector $1$-complete
condition. The theorem \path{actual_starred_input_certificate}
in \texttt{QLengthStarredCoordinateOrder} verifies the order of the
genuine marked-point coordinates returned by the source's
noncomputable coordinate map. These declarations passed
the Lean build and transitive axiom audit (PRs~147--148).
The arbitrary-height family above is an elementary extension;
it has not yet been formalized uniformly in $n$.

We would welcome clarification if an additional constraint on the
points in $W^{\mathcal D}_*$ or on the map $Q$ was intended but is not
stated in Definitions~20 and the proof of Lemma~27. Without such a
constraint, the tree component of the displayed $Q$ need not be skew.

\end{document}